\chapter{The kiosk}
\label{ch:kiosk}

\chapterbrief
  {A touchscreen that takes payments --- without a single line of crypto in the UI.}
  {The Adapter pattern, and threading discipline.}
  {Tap a drink, tap a card, watch it brew.}

\section{The problem}

You have a working payment stack. Now you need a screen. Two constraints make
this harder than it looks:

\begin{enumerate}
  \item You want to build the UI \emph{before} the hardware works --- or at least
        be able to develop it on a laptop with no reader attached.
  \item Card operations take \emph{seconds} (polling, RF, crypto), and a UI that
        blocks for seconds is a broken UI.
\end{enumerate}

\section{The idea: one interface, two implementations}

Define what the UI needs in the smallest possible terms, then satisfy it twice.

\begin{center}
\begin{tikzpicture}[font=\footnotesize\sffamily,
   b/.style={draw=brewernavy, rounded corners=2pt, inner sep=5pt,
             minimum width=42mm, minimum height=8mm},
   i/.style={draw=brewerblue, fill=brewerblue!8, rounded corners=2pt,
             inner sep=5pt, minimum width=42mm, minimum height=8mm}]
  \node[b] (ui) {JavaFX UI};
  \node[i, below=9mm of ui] (if) {\code{PaymentTerminal}};
  \node[b, below left=13mm and 6mm of if] (sim) {\code{SimulatedPaymentTerminal}};
  \node[b, below right=13mm and 6mm of if] (real) {\code{DesfirePaymentTerminal}};
  \draw[-{Stealth[length=4pt]}, brewernavy] (ui) --
       node[right, font=\scriptsize\itshape, color=brewergrey] {depends only on} (if);
  \draw[-{Stealth[length=4pt]}, dashed, brewergrey] (sim) -- (if);
  \draw[-{Stealth[length=4pt]}, dashed, brewergrey] (real) -- (if);
  \node[font=\scriptsize\itshape, color=brewergrey, below=1mm of sim] {laptop, no hardware};
  \node[font=\scriptsize\itshape, color=brewergrey, below=1mm of real, align=center]
       {your whole card stack};
\end{tikzpicture}
\end{center}

The same seam idea as \code{SpiLink} and \code{ApduChannel}, now at the top of
the stack. It means the entire UI can be built and tested before a single card
exists, and the kiosk degrades gracefully: no reader attached, so use the
simulator, and the app still runs.

\section{The contract}

\begin{jcode}
package com.midnightbrewer.payment;

/**
 * The seam between the kiosk UI and whatever actually takes the money.
 *
 * The UI depends only on this type, so it neither knows nor cares whether
 * payment is a simulated wallet or a real DESFire being debited over SPI.
 *
 * Extends AutoCloseable so the RF field is always powered down, even on a
 * crash.
 */
public interface PaymentTerminal extends AutoCloseable {

    /**
     * Begin waiting for a card and charge it for drink.
     *
     * MUST return immediately -- all polling happens on a background thread.
     * Exactly one terminal callback (approved or declined) is delivered per
     * call, unless cancel() intervenes.
     */
    void startTransaction(Drink drink, PaymentListener listener);

    /** Abort the in-flight transaction. Safe to call when nothing is running. */
    void cancel();

    /** Human-readable state for the status pill, e.g. "RC522 READY". */
    String getStatusText();

    @Override
    void close();
}
\end{jcode}

\begin{jcode}
public interface PaymentListener {
    /** A card entered the field. Not yet authenticated. */
    void onCardDetected(String uid);
    /** Authenticated and balance read, before any debit. */
    void onBalanceRead(int balanceCents);
    /** Debit COMMITTED. The drink may now be dispensed. */
    void onApproved(PaymentReceipt receipt);
    /** Failed. The card was not charged, or the debit was rolled back. */
    void onDeclined(CardException cause);
}
\end{jcode}

\begin{concept}{Why callbacks, and why four of them}
\textbf{Why not a return value?} Because \code{startTransaction} returns
instantly and the result arrives seconds later. A blocking \code{boolean
pay(Drink)} would freeze the UI thread: the animation stops, touches queue up,
and the operating system offers to kill your app.

\textbf{Why four callbacks?} So the UI can narrate. ``Card detected\ldots'' then
``Balance \texteuro5.00, charging\ldots'' then approved. A payment terminal that
shows nothing for four seconds feels broken even when it is working perfectly.
\end{concept}

\section{The Adapter}
\index{Adapter pattern}

Now the interesting class: something that satisfies \code{PaymentTerminal} by
driving the stack you built in Parts~II--IV.

\begin{analogy}{The travel plug adapter}
Your laptop charger has one plug shape; the Italian wall has another. The adapter
does not generate electricity or modify the laptop --- it sits between two fixed
interfaces and translates.

Both sides were designed without knowledge of each other, and the adapter makes
them fit. Here the kiosk speaks \code{PaymentTerminal} and \code{CardException};
your stack speaks \code{Rc522Reader}, \code{DesfireCard} and
\code{SpiException}. \code{DesfirePaymentTerminal} translates --- including the
exceptions, so a \code{com.midnightbrewer.hardware} type never leaks into the UI.
\end{analogy}

\begin{jcode}
package com.midnightbrewer.payment;

public class DesfirePaymentTerminal implements PaymentTerminal {

    // How the card was personalised (Chapter 12 / Appendix A).
    private static final byte[] APP_AID       = { 0x03, 0x02, 0x01 };
    private static final byte   USER_KEY_NO   = 0x02;
    private static final byte   VALUE_FILE_NO = 0x01;

    private final Pi4jSpiLink link;
    private final Rc522Driver driver;
    private final Rc522Reader reader;

    // ONE background thread, created once. Transactions are strictly
    // sequential -- two overlapping card conversations would interleave
    // command counters and corrupt both.
    private final ExecutorService worker =
            Executors.newSingleThreadExecutor(r -> {
                Thread t = new Thread(r, "desfire-terminal");
                t.setDaemon(true);      // never blocks JVM shutdown
                return t;
            });

    /** Written by the UI thread, read by the worker. */
    private final AtomicBoolean cancelled = new AtomicBoolean(false);
    private volatile boolean scanning = false;

    /** Opens the reader. THROWS if there is no RC522 -- see createTerminal(). */
    public DesfirePaymentTerminal() throws SpiException {
        this.link = new Pi4jSpiLink();
        this.driver = new Rc522Driver(link);
        try {
            initReader();
            int version = driver.version();
            if (version != 0x91 && version != 0x92) {
                throw new SpiException("RC522 not responding");
            }
        } catch (SpiException e) {
            quietClose();               // do not leak SPI on a failed construction
            throw e;
        }
        this.reader = new Rc522Reader(driver);
    }

    @Override
    public void startTransaction(Drink drink, PaymentListener listener) {
        cancelled.set(false);
        worker.submit(() -> runTransaction(drink, listener));   // returns instantly
    }

    private void runTransaction(Drink drink, PaymentListener listener) {
        scanning = true;
        try {
            // Fresh reader state each tap: softReset + antennaOn power-cycles
            // the field so the card is IDLE again and activate() runs cleanly
            // even after a previous purchase.
            initReader();

            if (!waitForCard()) return;         // cancelled; the UI moved on

            ActivatedCard card = reader.activate();
            listener.onCardDetected(card.uidHex());

            DesfireCard desfire = new DesfireCard(
                    new IsoDepChannel(reader), RandomSource.secure());
            desfire.selectApplication(APP_AID);
            desfire.authenticateEv2First(USER_KEY_NO, userKey());

            int balance = desfire.getValue(VALUE_FILE_NO);
            listener.onBalanceRead(balance);

            int price = drink.getPriceCents();
            if (balance < price) {
                listener.onDeclined(new InsufficientFundsException(balance, price));
                return;
            }

            // Provisional until the commit is acknowledged; only THEN may we
            // hand out a receipt, and only then may the UI dispense.
            desfire.debit(VALUE_FILE_NO, price);
            desfire.commitTransaction();
            int after = desfire.getValue(VALUE_FILE_NO);

            listener.onApproved(
                new PaymentReceipt(card.uidHex(), drink, balance, after));

        } catch (SpiException e) {
            listener.onDeclined(translate(e));  // stack vocabulary -> UI vocabulary
        } catch (RuntimeException e) {
            listener.onDeclined(new CardCommunicationException(
                    "unexpected reader error: " + e.getMessage(), e));
        } finally {
            scanning = false;
        }
    }

    private boolean waitForCard() throws SpiException {
        while (!cancelled.get()) {
            if (reader.isCardPresent()) return true;
            sleep(30);
        }
        return false;
    }
}
\end{jcode}

\begin{concept}{\texttt{catch (RuntimeException)} is not paranoia}
A kiosk must never die on a bad tap. \code{SpiException} covers expected
failures; the \code{RuntimeException} catch covers the unexpected --- an array
index, a null, a Pi4J internal error.

Without it, one odd card kills the worker thread and the machine stops taking
payments until someone reboots it.

\textbf{At a process boundary, catch broadly and report. Inside your logic, catch
narrowly.}
\end{concept}

\section{Threading: the rule you must not break}

JavaFX has exactly one thread allowed to touch the UI. Your terminal calls back
from a \emph{different} thread. So every callback hops:

\begin{jcode}
private final class UiPaymentListener implements PaymentListener {

    @Override
    public void onCardDetected(String uid) {
        Platform.runLater(() ->
            paymentHintLabel.setText("Card detected -- authenticating..."));
    }

    @Override
    public void onApproved(PaymentReceipt receipt) {
        Platform.runLater(() -> { stopPulse(); startBrewing(receipt); });
    }
}
\end{jcode}

\begin{trap}{Off-thread UI access does not reliably throw}
Sometimes it works. Sometimes it corrupts rendering in ways that appear once an
hour, on the real panel only, and never under a debugger.

This is one of the nastiest bug classes in UI programming, and the fix is a
habit, not a diagnosis: \textbf{every callback that touches the scene graph is
wrapped in \code{Platform.runLater}.} No exceptions.

The same rule exists everywhere: Android has \code{runOnUiThread}, Swing has
\code{invokeLater}, Qt has queued connections. One thread owns the widgets.
\end{trap}

\section{Choosing the backend}

The whole design collapses into one method:

\begin{jcode}
private PaymentTerminal createTerminal() {
    if ("simulated".equalsIgnoreCase(System.getProperty("brewer.reader", ""))) {
        return new SimulatedPaymentTerminal();          // forced, for UI work
    }
    try {
        PaymentTerminal real = new DesfirePaymentTerminal();
        System.out.println("[UI] RC522 ready -- real card payment enabled");
        return real;
    } catch (Exception e) {
        System.out.println("[UI] no RC522 (" + e.getMessage()
                + ") -- using the simulator");
        return new SimulatedPaymentTerminal();
    }
}
\end{jcode}

\textbf{Detection by attempting to open.} On a laptop, constructing
\code{Pi4jSpiLink} throws and you fall back. The same build runs everywhere: no
flags, no config file, no \code{\#ifdef}. This is precisely why the constructor
\emph{must} throw rather than returning a half-built object.

And because everything above depends on the interface, the UI needs
\textbf{zero} changes to switch backends. The bench ``simulate tap'' button even
hides itself:

\begin{jcode}
devTapButton.setVisible(terminal instanceof SimulatedPaymentTerminal);
\end{jcode}

\section{Verify}

\textbf{On a laptop} (no hardware): \code{mvn javafx:run -Dbrewer.windowed=true}.
The simulator's tap button drives the whole flow.

\textbf{On the Pi}, run a headless smoke test first --- it exercises the exact
path a tap takes, minus the screen:

\begin{jcode}
public final class TerminalCheck {
    public static void main(String[] args) throws Exception {
        Drink testDrink = new Drink("test", "Test Shot", 10);   // 10 cents
        try (DesfirePaymentTerminal terminal = new DesfirePaymentTerminal()) {
            CountDownLatch done = new CountDownLatch(1);
            terminal.startTransaction(testDrink, new PaymentListener() {
                public void onCardDetected(String uid)   { print("card: " + uid); }
                public void onBalanceRead(int cents)     { print("balance: " + cents); }
                public void onApproved(PaymentReceipt r) { print("APPROVED"); done.countDown(); }
                public void onDeclined(CardException e)  { print("DECLINED"); done.countDown(); }
            });
            System.out.println("Tap a card...");
            done.await();
        }
    }
}
\end{jcode}

\code{CountDownLatch} is how \code{main} waits for a background thread: start at
one, block on \code{await()}, and the callback counts down to zero to release it.

\begin{verifybox}
\begin{trace}
  card: 04 76 18 62 7B 20 90
  balance: $5.00
  APPROVED: charged $0.10  ($5.00 -> $4.90)
\end{trace}
\end{verifybox}

\section{Traps}

\begin{traptable}
UI freezes on tap & Card work on the FX thread. It must be on the worker. \\
Random rendering corruption & A callback touching nodes without \code{Platform.runLater}. \\
Second tap always fails & Reader not re-initialised per transaction. \\
Kiosk dies on an odd card & Missing the broad \code{RuntimeException} catch. \\
Simulator used on the Pi & Constructor threw --- read the printed message, it names the cause. \\
Drink dispensed but card not charged & Dispensing on \code{debit} instead of after \code{commitTransaction}. \\
\end{traptable}

\section{Commit}

\begin{shell}
git add -A
git commit -m "kiosk: PaymentTerminal seam + DesfirePaymentTerminal adapter"
\end{shell}
